\documentclass[10pt,a4paper]{amsart}
\usepackage{amsmath,amssymb,amsthm,graphicx}
% Retained minimal preamble for the Stacks Project chapter more-algebra.tex (More on Algebra).
% The chapter file pinned-original-source.tex carries only \input{preamble} on its first line, so no
% preamble text lives inside the chapter itself.  The Stacks Project uses one single shared preamble for
% every chapter; the snapshot stages it identically as sources/stacks-chow/preamble.tex and
% sources/stacks-etale-cohomology/preamble.tex.  The block below is copied mechanically, line for line,
% from sources/stacks-chow/preamble.tex lines 14-18 (the Xy-pic setup) and 154-217 (the native theorem
% declarations together with the mathematical macros).  Nothing was added, removed or redefined; the
% wrapper class and the amsmath/amssymb/amsthm/graphicx base packages are supplied by the assembler.
% Chapter: more-algebra.tex -- More on Algebra.
% For commutative diagrams we use Xy-pic
\usepackage[all]{xy}

% We use 2cell for 2-commutative diagrams.
\xyoption{2cell}

% Theorem environments.
%
\theoremstyle{plain}
\newtheorem{theorem}[subsection]{Theorem}
\newtheorem{proposition}[subsection]{Proposition}
\newtheorem{lemma}[subsection]{Lemma}

\theoremstyle{definition}
\newtheorem{definition}[subsection]{Definition}
\newtheorem{example}[subsection]{Example}
\newtheorem{exercise}[subsection]{Exercise}
\newtheorem{situation}[subsection]{Situation}

\theoremstyle{remark}
\newtheorem{remark}[subsection]{Remark}
\newtheorem{remarks}[subsection]{Remarks}

\numberwithin{equation}{subsection}

% Macros
%
\def\lim{\mathop{\mathrm{lim}}\nolimits}
\def\colim{\mathop{\mathrm{colim}}\nolimits}
\def\Spec{\mathop{\mathrm{Spec}}}
\def\Hom{\mathop{\mathrm{Hom}}\nolimits}
\def\Ext{\mathop{\mathrm{Ext}}\nolimits}
\def\SheafHom{\mathop{\mathcal{H}\!\mathit{om}}\nolimits}
\def\SheafExt{\mathop{\mathcal{E}\!\mathit{xt}}\nolimits}
\def\Sch{\mathit{Sch}}
\def\Mor{\mathop{\mathrm{Mor}}\nolimits}
\def\Ob{\mathop{\mathrm{Ob}}\nolimits}
\def\Sh{\mathop{\mathit{Sh}}\nolimits}
\def\NL{\mathop{N\!L}\nolimits}
\def\CH{\mathop{\mathrm{CH}}\nolimits}
\def\proetale{{pro\text{-}\acute{e}tale}}
\def\etale{{\acute{e}tale}}
\def\QCoh{\mathit{QCoh}}
\def\Ker{\mathop{\mathrm{Ker}}}
\def\Im{\mathop{\mathrm{Im}}}
\def\Coker{\mathop{\mathrm{Coker}}}
\def\Coim{\mathop{\mathrm{Coim}}}

% Boxtimes
%
\DeclareMathSymbol{\boxtimes}{\mathbin}{AMSa}{"02}

%
% Macros for moduli stacks/spaces
%
\def\QCohstack{\mathcal{QC}\!\mathit{oh}}
\def\Cohstack{\mathcal{C}\!\mathit{oh}}
\def\Spacesstack{\mathcal{S}\!\mathit{paces}}
\def\Quotfunctor{\mathrm{Quot}}
\def\Hilbfunctor{\mathrm{Hilb}}
\def\Curvesstack{\mathcal{C}\!\mathit{urves}}
\def\Polarizedstack{\mathcal{P}\!\mathit{olarized}}
\def\Complexesstack{\mathcal{C}\!\mathit{omplexes}}
% \Pic is the operator that assigns to X its picard group, usage \Pic(X)
% \Picardstack_{X/B} denotes the Picard stack of X over B
% \Picardfunctor_{X/B} denotes the Picard functor of X over B
\def\Pic{\mathop{\mathrm{Pic}}\nolimits}
\def\Picardstack{\mathcal{P}\!\mathit{ic}}
\def\Picardfunctor{\mathrm{Pic}}
\def\Deformationcategory{\mathcal{D}\!\mathit{ef}}
\usepackage{cleveref}
\crefname{originallemma}{original lemma}{original lemmas}
\Crefname{originallemma}{Original lemma}{Original lemmas}
\makeatletter
\newlabel{more-algebra-lemma-kunneth-converges}{{0H7W}{source}{}{sourcepointer-1}{}}
\newlabel{more-algebra-lemma-kunneth-converges@cref}{{[originallemma][0][]0H7W}{[1][0][]source}{}{}{}}
\makeatother
\crefname{originallemma}{original lemma}{original lemmas}
\Crefname{originallemma}{Original lemma}{Original lemmas}
\makeatletter
\newlabel{homology-lemma-biregular-ss-converges}{{012W}{source}{}{sourcepointer-2}{}}
\newlabel{homology-lemma-biregular-ss-converges@cref}{{[originallemma][0][]012W}{[1][0][]source}{}{}{}}
\makeatother
\crefname{originallemma}{original lemma}{original lemmas}
\Crefname{originallemma}{Original lemma}{Original lemmas}
\makeatletter
\newlabel{homology-lemma-ss-converges-trivial}{{0BK5}{source}{}{sourcepointer-3}{}}
\newlabel{homology-lemma-ss-converges-trivial@cref}{{[originallemma][0][]0BK5}{[1][0][]source}{}{}{}}
\makeatother
\title{Convergence of the Kunneth spectral sequences for a bounded complex}
\author{The Stacks Project Authors}
\date{stacks-project:more-algebra.tex:local-snapshot-20261012; GNU FDL 1.2 or later}
\begin{document}
\maketitle
\setcounter{subsection}{1}
\noindent\textit{Editorial note.} All delivered mathematical text is the author's own, copied from the pinned Stacks Project chapter file more-algebra.tex; no statement or proof step was rewritten, reordered or removed.\par\medskip
\noindent\textit{Editorial note.} The chapter file pinned-original-source.tex carries only \textbackslash{}input\{preamble\} on its first line, so no preamble text lives inside the chapter. The retained file sources/stacks-more-algebra/retained-preamble.tex is a mechanical, line-for-line copy of the project's single shared preamble as staged in the snapshot (Xy-pic setup plus theorem declarations and mathematical macros). Nothing was added, removed or redefined, and the excerpt serves only as preamble, never as a span coordinate.\par\medskip
\noindent\textit{Editorial note.} Both counted spans start at the \textbackslash{}begin line of their environment, so the statement and the author's proof are delivered whole; surrounding prose of the section is deliberately outside the counted spans and is carried as standard prerequisites instead.\par\medskip
\noindent\textit{Editorial note.} The citation fields (source\_id, source\_version, authors, source\_url, license.name, source\_title) are identical to those of every other card this lane has produced from this chapter, so the batch selector treats them as one source.\par\medskip
\noindent\textit{Allowed standard prerequisites (editorial).}\begin{itemize}
\item Spectral sequence of a filtered complex and the convergence criterion for a bounded filtration. used throughout the proof, lines 15543-15643. Supplies the object whose convergence is asserted and the criterion checked at the end.
\item The Kunneth spectral sequences for the tensor product and for the Hom complex. constructed in the same section, source lines 15300-15516, before the counted statement. Provides the two spectral sequences the statement refers to.
\item Identification of the graded pieces of a filtration with Tor and Ext terms. used at proof lines 15560-15620. Supplies the second page in both cases.
\end{itemize}\medskip
\ifdefined\hypertarget\hypertarget{sourcepointer-1}{}\fi
\noindent\textit{Source reference.} Original lemma 0H7W: focused-seeds/AGP/sources/stacks-snapshot-20261011/tags.txt, source lines 21471--21471. The lemma's own canonical tag, quoted so that readers can resolve the label inside the delivered text; it points at the counted statement itself and carries no proof content.\par\medskip
\ifdefined\hypertarget\hypertarget{sourcepointer-2}{}\fi
\noindent\textit{Source reference.} Original lemma 012W: focused-seeds/AGP/sources/stacks-snapshot-20261011/tags.txt, source lines 1414--1414. Navigation pointer to a statement the counted proof refers to but does not contain; nothing load-bearing is delegated to it.\par\medskip
\ifdefined\hypertarget\hypertarget{sourcepointer-3}{}\fi
\noindent\textit{Source reference.} Original lemma 0BK5: focused-seeds/AGP/sources/stacks-snapshot-20261011/tags.txt, source lines 14548--14548. Navigation pointer to a statement the counted proof refers to but does not contain; nothing load-bearing is delegated to it.\par\medskip
\begin{lemma}
\label{lemma-kunneth-converges}
Under the assumptions above, if in addition
\begin{enumerate}
\item the filtration on $K^\bullet$ is finite and the filtration on
$L^\bullet$ is finite, or
\item more generally the following are true
\begin{enumerate}
\item $F^iK^\bullet$ is acyclic for $i \gg 0$,
\item $F^iK^\bullet \to K^\bullet$ is a quasi-isomorphism for $i \ll 0$,
\item $F^jL^\bullet$ is acyclic for $j \gg 0$, and
\item $F^jL^\bullet \to L^\bullet$ is a quasi-isomorphism for $j \ll 0$.
\end{enumerate}
\end{enumerate}
Then the spectral sequence is bounded, the associated filtration on each
$H^n(T^\bullet) = H^n(K^\bullet \otimes_R^\mathbf{L} L^\bullet)$ is finite
and we have convergence
$$
E_1^{p, q} =
\bigoplus\nolimits_{i + j = p}
H^{p + q}(\text{gr}^iK^\bullet \otimes_R^\mathbf{L} \text{gr}^jL^\bullet)
\Rightarrow
H^n(K^\bullet \otimes_R^\mathbf{L} L^\bullet)
$$
\end{lemma}

\begin{proof}
In case (1) the filtration on $T^\bullet$ is finite and the lemma follows
immediately from
Homology, Lemma \ref{homology-lemma-biregular-ss-converges}.
In case (2) choose $a < b$ such that $F^iK^\bullet$ and $F^jL^\bullet$ are
acyclic for $i, j > b$ and $F^iK^\bullet \to K^\bullet$ and
$F^jL^\bullet \to L^\bullet$ are quasi-isomorphisms for $i, j < a$.
We claim that in this case the complex $F^nT^\bullet$ is acyclic for
$n > 2b$ and that $F^nT^\bullet \to T^\bullet$ is a quasi-isomorphism
for $n < 2a - 1$. Since the claim shows that
Homology, Lemma \ref{homology-lemma-ss-converges-trivial}
applies we see that our lemma is true.

\medskip\noindent
Proof of the claim. Given $n$ and integers $i_1 \leq i_2$ consider
$$
S^{n, i_1, i_2} =
\Im\left(
\bigoplus\nolimits_{i_1 \leq i \leq i_2}
\text{Tot}(F^iK^\bullet \otimes_R F^{n - i}L^\bullet)
\to
\text{Tot}(K^\bullet \otimes_R L^\bullet)
\right)
$$
viewed as a subcomplex of $\text{Tot}(K^\bullet \otimes_R L^\bullet)$.
Observe that $F^nT^\bullet = \colim S^{n, i_1, i_2}$ is a filtered colimit.
Thus to show the claim it suffices to show that $S^{n, i_1, i_2}$ is acyclic
for $n > 2b$ and that
$S^{n, i_1, i_2} \to \text{Tot}(K^\bullet \otimes L^\bullet)$
is a quasi-isomorphism for $n < 2a - 1$ and a cofinal set of choices of
pairs $i_1, i_2$.

\medskip\noindent
If $i_1 = i_2 = i$ and $n > 2b$, then
$$
S^{n, i, i} =
\text{Tot}(F^iK^\bullet \otimes_R F^{n - i}L^\bullet)
$$
is acyclic: either $i > b$ and this complex is acyclic by 
our assumption that $F^{n - i}L^\bullet$ is K-flat, or
$i \leq b$ and then $n - i \geq n - b > b$ and the same holds.
Using our flatness assumptions the reader shows that there is a short
exact sequence
$$
0 \to
S^{n + 1, i_2 + 1, i_2 + 1}
\to
S^{n, i_1, i_2} \oplus S^{n, i_2 + 1, i_2 + 1}
\to
S^{n, i_1, i_2 + 1}
\to
0
$$
of complexes. Thus we see by induction on $i_2 - i_1$ that all of the complexes
$S^{n, i_1, i_2}$ for $n > 2b$ are acyclic.

\medskip\noindent
Assume $n < 2a - 1$. Observe that
$$
S^{n, a - 1, a - 1} =
\text{Tot}(F^{a - 1}K^\bullet \otimes_R F^{n - a + 1}L^\bullet)
\to \text{Tot}(K^\bullet \otimes_R L^\bullet)
$$
is a quasi-isomorphism because both $a - 1 < a$ and $n - a + 1 < a$
and all our complexes are K-flat hence both sides compute the
same object of $D(R)$.
Observe that for $i_2 \geq a - 1$ and $n < 2a - 1$ the map
$$
S^{n + 1, i_2 + 1, i_2 + 1} \to S^{n, i_2 + 1, i_2 + 1}
$$
is a quasi-isomorphism because left and right hand side are quasi-isomorphic
to $\text{Tot}(F^{i_2 + 1}K^\bullet \otimes L^\bullet)$.
Thus using the short exact
sequence of complexes in the previous paragraph we find that
$S^{n, a - 1, a - 1} \to S^{n, a - 1, t}$ is a quasi-isomorphism for
all $t \geq a - 1$. On the other hand, there are similarly short
exact sequences
$$
0 \to
S^{n + 1, i_1, i_1}
\to
S^{n, i_1 - 1, i_1 - 1} \oplus
S^{n, i_1, i_2}
\to
S^{n, i_1 - 1, i_2}
\to
0
$$
of complexes. Similarly to the above, we note that for $i_1 \leq a - 1$
the map
$$
S^{n + 1, i_1, i_1} \to S^{n, i_1 - 1, i_1 - 1}
$$
is a quasi-isomorphism because left and right hand side are quasi-isomorphic
to $\text{Tot}(K^\bullet \otimes F^{n - i_1 + 1}L^\bullet)$.
Thus we conclude that $S^{n, s, t} \to S^{n, a - 1, t}$ is a quasi-isomorphism
for all $s \leq a - 1$. Combined we find that $S^{n, s, t}$ for $s \leq a - 1$
and $t \geq a - 1$ maps quasi-isomorphically to
$\text{Tot}(K^\bullet \otimes_R L^\bullet)$.
This finishes the proof of the lemma.
\end{proof}

\end{document}
